\chapter{What this toolkit reaches}

\textbf{Purpose:} Record the recommendation made about which problem the project's apparatus could
be pointed at, the reasoning offered for it, the trap identified alongside it, and the measured
bound that constrains it. \textbf{Scope:}
\texttt{BTCminer/\allowbreak{}examples/\allowbreak{}proofing/\allowbreak{}bbp\_\allowbreak{}scaling.\allowbreak{}py},
\texttt{BTCminer/\allowbreak{}theory/\allowbreak{}PARTITION.\allowbreak{}tsv}.

\textbf{Whose this is.} The recommendation in this chapter is the project architect's, given on
2026-09-11 in answer to a direct question, and is recorded for Douglas Quigg to rule on. It is not
this book's recommendation and no work has begun on it.

\section{The trap, which is the obvious answer}

Riemann is what an engine whose edge is arbitrary precision reaches for first, and it is the wrong
choice for a reason that has nothing to do with how good the arithmetic is.

The apparatus matches the \emph{subject} and not the \emph{obstruction}. A zeta function and
unbounded precision look like a fit. But computing zeros on the critical line is long done, far past
anything a desktop reaches, and no finite count of zeros bears on the proof at all. Adding zeros
reproduces existing work at a worse scale. The conjecture is not precision-bound, so an engine
whose advantage is precision buys nothing against it.

That is the superficial-match failure in its exact form, and it is worth having named because it is
the shape of mistake this apparatus is most exposed to: the subject looks like the instrument.
Yang--Mills wants constructive quantum field theory and lattice gauge machinery. Hodge wants
algebraic geometry with no numerical handle this toolkit has. P versus NP is not a numerical
question, and Douglas Quigg had already set it aside on the separate ground that the anisotropic
boundary has not been worked out.

\textbf{The specific figures on Riemann given for this argument were not independently checked.}
The claim that zeros have been computed at height $10^{22}$ and that verification has covered the
first $10^{13}$ of them was supplied with the recommendation and is recorded as supplied. It is not
load-bearing for the argument, which turns on the conjecture not being precision-bound rather than
on any particular height reached.

\section{The recommendation}

Birch and Swinnerton-Dyer, on the ground that it is the one Millennium problem whose obstruction
matches the instrument rather than whose subject does.

It is the only one of them found by computation: Birch and Swinnerton-Dyer fitting curves
numerically on EDSAC. It is natively experimental, and the refined conjecture is an exact numerical
identity. For an elliptic curve $E$ of rank $r$, writing $\mathrm{Sha}$ for the order of the
Tate--Shafarevich group, $\Omega$ for the real period, $\mathrm{Reg}$ for the regulator and $c_p$
for the Tamagawa numbers,
\[
  \frac{L^{(r)}(E,1)}{r!} \;=\; \frac{\mathrm{Sha}\cdot\Omega\cdot\mathrm{Reg}\cdot\prod_p c_p}
                                    {|E_{\mathrm{tors}}|^2}.
\]
Rearranged, the analytic order of $\mathrm{Sha}$ is a quotient of computable real numbers, and the
conjecture asserts that this quotient is an integer, usually a perfect square. The verification is
therefore: compute a real number to $N$ digits and ask whether it lands on an integer. That is a
precision-bound question, and it is the check the engine already performs.

\section{The first measurement, and the positive control it requires}

The positive control comes first and is not optional. Curves are taken from published tables where
the rank and the order of $\mathrm{Sha}$ are already known, at rank zero and one, small conductor,
with $\mathrm{Sha}$ trivial or known to be 4 or 9. The quotient is computed and must land on the
known integer. An instrument that cannot reproduce a known answer says nothing when pointed at an
unknown one.

\textbf{The quotient is only as good as its weakest input.} The real period and the regulator come
from their own algorithms carrying their own error. The honest report is therefore the precision of
the quotient and not the precision of the arithmetic. An arithmetic with no floor under it does not
give its inputs one.

\section{The null: the derived version is withdrawn}

A null was offered with the recommendation and then withdrawn by its author in the same exchange.
The withdrawal is the substantive part and is recorded here rather than only in the ledger.

\textbf{Withdrawn.} That a real number carrying no integrality structure has its distance to the
nearest integer distributed uniformly on $[0,\tfrac12]$, so that landing within $10^{-N}$ of an
integer by chance has probability about $2\times10^{-N}$, and agreement to fifty digits is
therefore a bound of $10^{-50}$.

\textbf{What killed it.} It was derived and not drawn. A threshold reasoned out from an assumed
distribution omits the variance sources its author failed to enumerate, and omitted variance only
ever pushes one way. The rule in this project is to draw the bar and never to derive it, and this
null was derived in the same message that stated the positive-control rule.

One narrower objection raised against it was itself wrong and is recorded as wrong. The suggestion
was that the integer factors in the quotient, the torsion order squared, the Tamagawa product and
$r!$, might concentrate the quotient near integers and inflate the coincidence rate. They do not:
scaling a structureless real by a known rational does not concentrate it near integers. The null was
wrong for the better reason and not for that one.

\textbf{The drawn replacement} measures the null under the same conditions as the effect, through
the same pipeline and with the same inputs and the same precision losses, in two ways. Forming the
quotient with the rank taken one too high or one too low gives a condition under which the identity
should not hold, and the distribution of distance-to-nearest-integer across many curves is then the
null itself. Computing the quotient across many curves whose $\mathrm{Sha}$ is known gives the
spread of the residual about the true integer, which is the instrument's own floor and which will be
dominated by the period and the regulator rather than by the arithmetic. The second is the positive
control and the floor measurement in one pass.

The consequence changes what may be claimed. The honest bound is not $10^{-N}$ for $N$ digits of
arithmetic. It is however close the pipeline gets on a curve whose answer is already known. If the
regulator runs out at $10^{-12}$, then $10^{-12}$ is the claim however many digits the multiply
carries.

\section{The measured bound that constrains this, and it was checked}

\texttt{bbp\_scaling.py} in the BTCminer tree measures what depth costs in the integer relation
search, and its recorded results were read in the file rather than taken on report.

The relation search costs precision raised to 2.64, measured, against 2.585 predicted from a
Karatsuba host multiply near width$^{1.585}$ times a count linear in the width. The reduction is not
the culprit. A device multiply only beats the host above 1{,}024 limbs, which is 32{,}768 bits, and
because the mean operand is 0.642 times the precision, the mean multiply reaches that crossover only
at about 51{,}000 binary places, where the file records one cell at 51{,}036 places costing 26.2
hours.

For the recommended measurement this is permissive: the precision that binds is tens to hundreds of
digits, nowhere near that wall. It bounds two other things. Nothing should be planned that needs
relation detection at tens of thousands of digits, and the card should not be expected to rescue it.
The file's own note records that the opposite was predicted before the measurement was taken.

\section{A precondition that has to be settled first}

The work spans two repositories with different governance, and this was established by checking
rather than assumed.

The arithmetic the recommendation depends on lives in the BTCminer tree and is classified EXTERNAL
in \texttt{PARTITION.\allowbreak{}tsv} under a ruling quoted there from Douglas Quigg on
2026-09-11, that advantage research is held and precision work is fair. The relation search, the
transform multiply and the arbitrary-precision arithmetic all carry that class. So the recommended
line of attack is not blocked by the partition.

\textbf{The adjacent material is blocked, and the boundary is narrow.} The same file classifies the
harmonic machinery as internal, under a separate ruling quoted there that anything used to derive
anisotropic properties is internal. A line of attack that begins wanting that machinery has crossed
into held material and stops for a ruling rather than proceeding.

The theory books are in a different repository from the arithmetic, with its own custodian. Anything
that wants files from both needs that settled before it starts and not after.

\section{What is not claimed}

This is verification and not proof, and the distinction is the whole of the problem.

A measurement of this kind produces evidence for particular instances of the conjecture, in the way
the original EDSAC work did. A verified instance at unusual rank or unusual precision is a real
contribution and it is not a solution.

The architect recording the recommendation stated plainly that they are not an expert in
elliptic-curve $L$-function computation, that computing $L^{(r)}(E,1)$ requires modular symbols or
Dokchitser's algorithm, that established systems already implement those and rebuilding them would
be foolish, and that the realistic shape is established algorithms for the $L$-function with this
engine supplying precision underneath. The first honest task is finding out whether that seam works
at all. That assessment is recorded as theirs and is not softened here.
